\documentclass{situationist-preprint}

\title{Creeping Minimalism}
\msid{MTE-17-b}
\msfamily{MTE-17}
\msdiscipline{Metaethics}
\mstype{Companion manuscript}
\mskeywords{creeping minimalism; quasi-realism; expressivism; explanation; deflationism; realism}

\begin{document}
\maketitle

\begin{abstract}
Dreier's problem of creeping minimalism is that, once the expressivist earns the right to talk of moral truth, facts and properties in a minimalist sense, there seems to be nothing left to distinguish her from the realist. Dreier proposes that the difference lies in explanation: the expressivist explains moral judgment without appeal to moral properties. We argue that creeping minimalism is the first symptom of the collapse established in MTE-17-a. The explanation criterion does distinguish expressivism, but as a theory of meaning, not as a view of what moral facts are; once the quasi-realist's account of error is made explicit, the explanation she offers of correctness is the realist's account of wrongness.
\end{abstract}

\section{The problem and Dreier's criterion}

Quasi-realism sets out to earn, on expressivist foundations, the realist's ways of talking: that moral claims are true, that there are moral facts and properties, that they are mind-independent \citep{Blackburn1993}. With a minimalist theory of truth, facts and properties, each of these is earned cheaply: \enquote{it is true that lying is wrong} says no more than that lying is wrong, and wrongness is a property in the sense that \enquote{wrong} is a predicate. \citet{Dreier2004} observed that, as minimalism creeps from truth to facts to properties to belief, every thesis the realist could use to state her view is one the quasi-realist also accepts, and the views threaten to become indistinguishable. His proposal was that the difference lies in explanation. The expressivist explains what it is to make a moral judgment, and what moral sentences mean, in terms of attitudes, without using moral vocabulary or appealing to moral properties; the realist explains them by the relation of judgments to the moral properties they represent.

\section{Where the criterion leads}

The explanation criterion is correct as far as it goes. MTE-03 shows that the semantics of moral sentences can be given in terms of states of being for, which compose according to the logic of plans, and this semantics uses no moral properties. In the theory of meaning, the expressivist's explanation differs from the representationalist's, and the difference survives minimalism.

The criterion does not, however, keep quasi-realism distinct from realism as a view of what moral facts are. The quasi-realist must also explain the realist-sounding things she earns, and in particular moral error: that a community, or the whole of humanity, could be wrong about what is wrong. Her explanation is that a judgment is mistaken when an improved sensibility would not share it \citep{Blackburn1998}. This is commitment (Q1) of MTE-17-a. If, as Dreier's criterion requires, the explanation is to proceed without appeal to moral properties, then the standard of improvement must be specified in non-moral terms, which is commitment (Q2). Under the deflationism that powers the creep, the account of when a moral judgment is correct is, by the equivalence schema, an account of when its content holds. Theorem~4.1 of MTE-17-a shows that the two commitments, given the equivalence schema and the coarse individuation of minimal properties, are equivalent to the realism of MTE-01, on which wrongness is a natural property of being disapproved under idealized conditions.

So the creep does not stop at the theory of meaning. Minimalism about truth turns the quasi-realist's account of correctness into an account of the property, and the explanation criterion, applied to that account, requires that the property be natural. The quasi-realist who satisfies Dreier's criterion everywhere is, in what she says about the nature of wrongness, a naturalist realist. If she relaxes the criterion and specifies improvement in moral terms, Corollary~4.3 of MTE-17-a shows that she becomes a non-naturalist realist, and the criterion fails for the account of error.

\section{Objection and reply}

\begin{objection}
The quasi-realist's account of error is not an account of what makes a moral judgment true but an expression of her own higher-order attitude toward sensibilities. It is therefore not an explanation in the sense the collapse requires.
\end{objection}

\begin{reply}
If the account of error merely expresses an attitude, it does not explain why we are entitled to say that we may be wrong, which was the quasi-realist's undertaking; it repeats the claim in a new key. Moreover, the expressed attitude has a content that can itself be assessed for correctness, and the quasi-realist owes an account of when that content is correct. The regress ends only with a standard of improvement, and the dilemma of MTE-17-a then applies to that standard \citep{Egan2007}.
\end{reply}

\section{Conclusion}

Creeping minimalism is real, and Dreier's explanation criterion is the right response to it in the theory of meaning, where the expressivist semantics of MTE-03 succeeds. But the creep continues past meaning. Once the quasi-realist explains error, the explanation, made explicit and freed of moral vocabulary, is the naturalist realist's account of wrongness. Expressivism's success, made explicit, collapses into realism.

\begin{thebibliography}{9}\small
\bibitem[Blackburn(1993)]{Blackburn1993} Blackburn, S. (1993). \emph{Essays in Quasi-Realism}. Oxford University Press.
\bibitem[Blackburn(1998)]{Blackburn1998} Blackburn, S. (1998). \emph{Ruling Passions: A Theory of Practical Reasoning}. Clarendon Press.
\bibitem[Dreier(2004)]{Dreier2004} Dreier, J. (2004). Meta-ethics and the problem of creeping minimalism. \emph{Philosophical Perspectives}, 18, 23--44.
\bibitem[Egan(2007)]{Egan2007} Egan, A. (2007). Quasi-realism and fundamental moral error. \emph{Australasian Journal of Philosophy}, 85(2), 205--219.
\end{thebibliography}

\end{document}
